\section{Анализ предметной области}

\subsection{Понятие и принципы онлайн-аренды вещей}

Электронная коммерция является одним из направлений цифровой экономики и охватывает совокупность торговых и финансовых операций, осуществляемых посредством компьютерных сетей, а также сопутствующие им бизнес-процессы. Онлайн-аренда вещей представляет собой отдельный сегмент электронной коммерции, в рамках которого пользователям предоставляется возможность временного использования материальных объектов через интернет-платформы на возмездной основе.

К основным моделям электронной коммерции относятся:

\begin{itemize}
\item P2P (peer-to-peer) --- взаимодействие между пользователями;
\item B2B (business-to-business) --- взаимодействие между организациями;
\item B2C (business-to-consumer) --- взаимодействие между компанией и конечным потребителем;
\item C2C (consumer-to-consumer) --- взаимодействие между частными потребителями;
\item B2G (business-to-government) --- взаимодействие между коммерческими организациями и государственными структурами.
\end{itemize}

Онлайн-бронирование представляет собой процесс резервирования товаров или услуг через сеть Интернет в режиме реального времени. Данный подход широко применяется при бронировании гостиничных номеров, приобретении билетов, резервировании мест в ресторанах, театрах и других сферах обслуживания.

Использование онлайн-платформ аренды обеспечивает ряд преимуществ для всех участников сделки:

\begin{enumerate}
    \item Экономическая эффективность. Пользователи получают возможность пользоваться дорогостоящими или редко используемыми вещами без необходимости их приобретения, тогда как владельцы получают дополнительный доход от временно неиспользуемого имущества.
    \item Простота доступа к услугам. Поиск объектов аренды, оформление бронирования и проведение оплаты доступны круглосуточно независимо от местоположения пользователя при наличии подключения к сети Интернет.
    \item Прозрачность взаимодействия. Система рейтингов и отзывов позволяет формировать репутацию участников платформы и способствует повышению доверия между сторонами сделки.
    \item Широкий выбор предложений. Пользователи получают возможность быстро сравнивать различные варианты аренды и выбирать наиболее подходящие условия.
\end{enumerate}

Несмотря на преимущества, онлайн-аренда сопровождается рядом рисков и ограничений:

\begin{enumerate}
    \item Отсутствие предварительного осмотра имущества. До момента получения вещи пользователь не может объективно оценить её состояние, вследствие чего возможны расхождения между фактическими характеристиками и информацией, указанной в объявлении.
    \item Логистические сложности и риски повреждения имущества. Передача объекта аренды может сопровождаться дополнительными затратами, а также рисками утраты или порчи имущества.
    \item Проблемы контроля исполнения обязательств. Возможны случаи несвоевременного возврата вещей, нарушения условий аренды или возникновения спорных ситуаций между сторонами.
    \item Юридические и финансовые споры. Возмещение ущерба, возврат залоговых сумм и урегулирование конфликтов могут требовать дополнительных механизмов правовой защиты.
    \item Угрозы информационной безопасности. Онлайн-платформы подвержены рискам мошенничества, несанкционированного доступа к персональным данным и финансовым операциям.
\end{enumerate}

\subsection{Классификация онлайн-площадок для аренды}

Современный рынок сервисов аренды характеризуется значительным разнообразием платформ, которые могут классифицироваться по нескольким признакам: ассортименту предлагаемых товаров, модели взаимодействия участников и типу арендодателей. Рассмотрение существующих классификаций позволяет определить особенности разрабатываемой системы и её целевую аудиторию.

\subsubsection{По типу предлагаемых товаров}

С точки зрения специализации платформы подразделяются на универсальные и специализированные. Универсальные сервисы предоставляют возможность аренды различных категорий товаров: электроники, строительного инструмента, спортивного инвентаря, одежды и других объектов. Их преимуществом является широкий выбор предложений и большая пользовательская база. Однако реализация специализированного функционала для каждой категории товаров в подобных системах затруднена.

Специализированные площадки ориентированы на ограниченный перечень товаров или одну конкретную категорию. Такие решения позволяют глубже учитывать особенности предметной области, внедрять специализированные механизмы взаимодействия и формировать сообщество пользователей с высоким уровнем экспертности.

\subsubsection{По модели взаимодействия с пользователями}

В зависимости от степени участия платформы в процессе заключения сделки можно выделить несколько моделей функционирования сервисов аренды.

Первая модель представляет собой классифайд-площадку, где сервис выполняет роль посредника для публикации объявлений, а участники самостоятельно договариваются о способах оплаты и условиях передачи имущества.

Вторая модель предполагает полное сопровождение сделки платформой. В данном случае система обеспечивает обработку платежей, хранение денежных средств до завершения сделки, разрешение конфликтных ситуаций и дополнительные механизмы защиты участников. Именно данная модель положена в основу проектируемого веб-приложения.

Третья модель основана на подписке и предусматривает регулярную оплату пользователем доступа к определённому набору объектов аренды. Наибольшее распространение такой подход получил в сфере одежды и детских товаров.

\subsubsection{По типу арендодателей}

По составу участников сделки онлайн-площадки подразделяются на две основные категории.

Модель C2C (Consumer-to-Consumer) предполагает взаимодействие между частными лицами, когда владельцы временно передают свои вещи другим пользователям за установленную плату.

Модель B2C (Business-to-Consumer) ориентирована на сотрудничество коммерческих организаций с конечными потребителями. В этом случае арендодателями выступают компании, профессионально предоставляющие услуги проката имущества.

\subsection{Основные бизнес-процессы системы аренды}

\subsubsection{Регистрация и верификация пользователей}

Для обеспечения безопасности сделок и снижения вероятности мошеннических действий платформа должна использовать многоуровневую систему регистрации пользователей. Каждый участник системы может выполнять функции арендатора либо арендодателя. Получение полного доступа к функционалу сервиса предусматривает обязательную процедуру подтверждения личности, включающую подтверждение адреса электронной почты и, при необходимости, предоставление дополнительных идентификационных данных через государственные сервисы или банковские инструменты.

\subsubsection{Публикация и поиск предложений}

Для размещения объекта аренды пользователь формирует карточку объявления, содержащую наименование вещи, категорию, описание, фотографии, стоимость аренды и периоды доступности. Со своей стороны платформа предоставляет средства структурирования информации, поиска и фильтрации объявлений по различным параметрам.

\subsubsection{Бронирование и совершение сделки}

Процесс оформления аренды автоматизирован за счёт использования встроенного механизма бронирования и календаря занятости. Такой подход исключает возможность одновременного оформления нескольких заказов на один объект аренды. После выбора периода аренды система автоматически рассчитывает итоговую стоимость услуги. Проведение платежей осуществляется через платформу, которая удерживает средства до подтверждения выполнения обязательств обеими сторонами сделки.

\subsubsection{Система рейтингов и отзывов}

Для сервисов, работающих по модели C2C, важную роль играет механизм формирования репутации пользователей. После завершения аренды участники могут оставлять отзывы и выставлять оценки друг другу. Накопленный рейтинг влияет на уровень доверия пользователей и ранжирование объявлений в результатах поиска.

\subsection{Компоненты программно-информационной системы}

Информационная система онлайн-площадки аренды представляет собой совокупность взаимосвязанных программных компонентов, обеспечивающих хранение, обработку и предоставление данных участникам платформы.

\subsubsection{Модуль управления пользователями}

Данный модуль реализует регистрацию, аутентификацию и авторизацию пользователей. Кроме того, он обеспечивает хранение профилей, сведений о совершённых сделках, рейтингах, отзывах и данных, используемых при прохождении процедуры верификации.

\subsubsection{Модуль каталога и поиска}

Модуль отвечает за создание, публикацию и редактирование объявлений, а также предоставляет механизмы полнотекстового поиска и фильтрации по категориям товаров и доступным датам аренды.

\subsubsection{Модуль бронирования и календаря}

Основной задачей данного компонента является контроль доступности объектов аренды и предотвращение конфликтов бронирования. Модуль поддерживает календарь занятости каждого объявления и обеспечивает уведомление участников о создании, подтверждении либо отклонении заявок.

\subsubsection{Модуль генерации договоров}

Компонент предназначен для автоматического формирования договоров аренды на основе информации из пользовательских профилей и карточек объявлений. Модуль обеспечивает создание документов, их подписание средствами электронной подписи, а также сопровождение процессов изменения и прекращения действия договоров.

\subsubsection{Модуль коммуникации}

Средства коммуникации реализованы посредством встроенного мессенджера, позволяющего участникам сделки обмениваться сообщениями без раскрытия персональных контактных данных. Дополнительно используются push-уведомления и электронные рассылки для информирования пользователей о значимых событиях системы.

\subsubsection{Модуль аналитики и отчётности}

Модуль предоставляет владельцам объявлений инструменты анализа эффективности использования платформы, включая сведения о доходах, популярности объектов аренды и количестве заключённых сделок. Для арендаторов формируются данные об истории бронирований и активности использования сервиса.

\subsection{Правовое регулирование сферы онлайн-аренды}

Деятельность онлайн-площадки по сдаче вещей в аренду регулируется рядом нормативно-правовых актов Российской Федерации. Основой является Гражданский кодекс РФ, а именно глава 34 «Аренда», которая определяет общие положения о договоре аренды, права и обязанности арендодателя и арендатора, ответственность за недостатки сданного в аренду имущества\cite{zakon2}.

Поскольку арендатор использует вещь для личных, семейных или домашних нужд, на отношения между участниками сделки распространяются положения Закона РФ «О защите прав потребителей» от 07.02.1992 № 2300-1. Согласно данному закону, арендатор имеет право на получение полной и достоверной информации об услуге, на отказ от исполнения договора в любое время при условии оплаты арендодателю фактически понесённых расходов, а также на возмещение убытков, причинённых вследствие недостатков оказанной услуги. Арендодатель, в свою очередь, обязан предоставить вещь в состоянии, соответствующем условиям договора аренды и назначению имущества, и несёт ответственность за недостатки, обнаруженные в течение срока аренды\cite{zakon1}.

Платформа, являясь оператором персональных данных, обязана соблюдать требования Федерального закона № 152-ФЗ «О персональных данных» от 27.07.2006 г. Вся собираемая информация о пользователях (паспортные данные, контактная информация, платёжные реквизиты) должна обрабатываться на законном основании, а также надёжно храниться и защищаться от несанкционированного доступа\cite{zakon4}. Для реализации функций безопасной сделки платформа интегрируется с платёжными сервисами, сертифицированными по стандарту PCI DSS или ISO-20022, что обеспечивает безопасность финансовых транзакций.

Важным аспектом правового регулирования является налогообложение доходов арендодателей. Арендодатели — физические лица, не зарегистрированные как индивидуальные предприниматели, могут применять специальный налоговый режим «Налог на профессиональный доход» в соответствии с Федеральным законом № 422-ФЗ от 27.11.2018 г\cite{zakon3}.

\subsection{Анализ существующих решений и аналогов}
Для улучшения качества собственного продукта следует провести анализ уже существующих решений. Сфера онлайн-бронирования с каждым годом развивается и расширяется. Наиболее востребованными в данном сегменте являются платформы Rentomania, Авито и Арендую.ру. Сравнение функционала данных площадок представлено в таблице \ref{concurent:table}.

\begin{xltabular}{\textwidth}{|X|X|X|X|X|}
 \caption{Сравнение существующих решений и аналогов\label{concurent:table}}\\ \hline
 \centrow Параметр & \centrow Rentomania & \centrow Авито & \centrow Арендую.ру & \centrow Проектируемая онлайн площадка \\ \hline
 \endfirsthead
 \continuecaption{Продолжение таблицы \ref{concurent:table}}
 ~ \centrow Параметр & \centrow Rentomania & \centrow Авито & \centrow Арендую.ру & \centrow Проектируемая онлайн площадка \\ \hline
 \finishhead
 \leftrow Размещение объявлений арендодателем & \centrow + & \centrow + & \centrow + & \centrow + \\ \hline
 \leftrow Поиск и фильтрация товаров & \centrow + & \centrow + & \centrow + & \centrow + \\ \hline
 \leftrow Онлайн-бронирование & \centrow + & \centrow -- & \centrow -- & \centrow + \\ \hline
 \leftrow Встроенная онлайн-оплата & \centrow + & \centrow + & \centrow -- & \centrow + \\ \hline
 \leftrow Система рейтинга и отзывов & \centrow + & \centrow + & \centrow -- & \centrow + \\ \hline
 \leftrow Чат между арендатором и арендодателем & \centrow + & \centrow + & \centrow -- & \centrow + \\ \hline
 \leftrow Адаптивная вёрстка для мобильных устройств & \centrow + & \centrow + & \centrow -- & \centrow + \\ \hline
 \leftrow Верификация пользователей & \centrow + & \centrow -- & \centrow -- & \centrow + \\ \hline
 \leftrow Панель аналитики для арендодателя & \centrow -- & \centrow -- & \centrow -- & \centrow + \\ \hline
 \leftrow Понятный интерфейс & \centrow -- & \centrow + & \centrow -- & \centrow + \\ \hline
 \leftrow Уведомления о статусе аренды & \centrow + & \centrow -- & \centrow -- & \centrow + \\ \hline
\end{xltabular}